\subsection{The height-six CM Jacobian has the required integral normal form}

\begin{proposition}[DP45: an actual three-direction CM example]\label{DP45}
Let $p\equiv3$ or $5\bmod7$, $R=W(\F_{p^6})$, and
$J=\operatorname{Jac}(y^2=x^7-1)$ with the good integral model of DP22.
Then $\widehat J$ is the actual connected height-six, dimension-three
$p$-divisible group with a unit-arrow cyclic display of type
$A=\{0,1,2\}\subset\Z/6\Z$.
In primitive integral eigendirections its sharp slopes are
\[
 (\gamma_0,\gamma_1,\gamma_2)
       =\frac{(p,p^2,p^3)}{(p-1)(p^3+1)}.
\]
Its actual dual has the complementary type and the same triple.
The mixed bound and exact $q^r$ coordinate witnesses of DP40 apply to this Jacobian, and DP43 gives optimal affine loss three in every direction.
The mean of its three slopes is
\[
 c_2\frac{p^2+p+1}{3(p^2-p+1)}<c_2,
\]
although the largest slope exceeds $c_2$ by $p/(p^3+1)$.
\end{proposition}

\begin{proof}
Since $p\ne2,7$, the smooth projective two-chart model
$y^2=x^7-1$, $v^2=u-u^8$ of DP22 has a genus-three Jacobian abelian scheme; the same FGA inputs cited there apply.
Let $a$ be the \emph{covariant} Jacobian automorphism
$a[P-\infty]=[(\zeta_7x(P),y(P))-\infty]$.
The integral basis $x^{i-1}dx/y$, $i=1,2,3$, has eigenvalues
$\zeta_7^i$ and orders $4,2,0$ at infinity. The induced cotangent action, and hence its ordinary transpose on $\Lie J$, has those same three characters. As in DP22, $\Phi_7(a)=0$: its differential on the characteristic-zero generic fibre vanishes, so the homomorphism vanishes there and hence on the abelian scheme.
The action extends to $\Z_p[\zeta_7]$ on $J[p^\infty]$.
Because $p$ has order six modulo seven, its fraction field is the unramified degree-six field $E_p$.

Write $G=J[p^\infty]$. Tate's connected--\'etale sequence is functorial.
If the \'etale quotient were nonzero, the unital $E_p$ action on its isogeny endomorphisms would be injective. Its rational Tate space would have dimension a positive multiple of six. Since $G$ has height six, the quotient would have full height and $G$ would be \'etale, contradicting its dimension three.
Thus $G$ is connected. The dual has the same height and, by the dimension formula, dimension three; the same field-action argument proves it connected. Tate's formal-group equivalence and abelian-scheme example identify $G$ with $\widehat J$.
These facts use \cite[\S1.4, p.~160; \S2.1 Example (a), p.~161; \S2.2 Proposition 1 and (4), pp.~162,164; \S2.3 Proposition 3 and Example (b), p.~166]{Tate67}.

Work first over $R_n=R/p^n$. Zink's Theorem~9, p.~132, supplies the display of this connected group over the excellent local $p$-nilpotent base; Theorem~5 gives its Witt rank six. Full faithfulness transports the cyclotomic action.
Let $\widehat\zeta=[\zeta_7]\in W(R_n)$. Since seven is invertible, the operators
\[
 e_b=\frac17\sum_{r=0}^6\widehat\zeta^{-br}a^r
       \quad(0\le b\le6)
\]
are orthogonal idempotents, with $e_0=0$ by $\Phi_7(a)=0$.
They split $P_n$ into six finite projective summands $P_{b,n}$.
The ring $W(R_n)$ is local (a Witt vector is a unit exactly when its zeroth component is a unit), so these summands are free.
Semilinear Frobenius sends $P_{b,n}$ to $P_{pb,n}$: it sends
$\widehat\zeta^b$ to $\widehat\zeta^{pb}$.
Its linearization is invertible after $p$ is inverted, since relative to any normal decomposition it differs from the invertible normal map by multiplying each $L$-column by $p$.
The localized Witt ring is nonzero, as it maps to
$W(\F_{p^6})[1/p]$.
The six free ranks are therefore equal along the single Frobenius cycle and sum to six. Each rank is one.

The Hodge sequence
\[
 0\longrightarrow Q_n/I_nP_n\longrightarrow P_n/I_nP_n
       \longrightarrow P_n/Q_n=\Lie G_{R_n}\longrightarrow0
\]
is equivariant. On characters $1,2,3$ its middle and right terms are free of rank one; their surjection is an isomorphism. On the other three characters the right term is zero. Hence
\[
 Q_n\cap P_{b,n}=
 \begin{cases}I_nP_{b,n},&b=1,2,3,\\P_{b,n},&b=4,5,6.\end{cases}
\]
Order the characters as
\[
 (1,3,2,6,4,5)\quad(p\equiv3\bmod7),\qquad
 (2,3,1,5,4,6)\quad(p\equiv5\bmod7).
\]
The first three positions are the tangent type. In chosen free generators, $F$ and $V^{-1}$ on the normal summands move to the next line. Their normal matrix is monomial and invertible, so its coefficients are units. This is exactly DP40's arbitrary-unit display.

The choices can be made compatibly in $n$: display base change preserves the idempotents, and a free-line generator lifts through the surjective Witt reductions to a generator. The resulting arrow units lie in
$\varprojlim_n W(R_n)=W(R)$.
Thus the formal isomorphism is integral at all levels, with primitive tangent basis; no isogeny rescaling is used.
DP40 identifies its actual dual with the complementary reciprocal-unit display. For clarity, $a^D$ acts on a covector by precomposition, retaining the eigenvalue $\widehat\zeta^b$, not inverting it. Its tangent characters are therefore the complementary three, consistent with the dual Hodge pairing in Zink Definition~19/(22).

The three numerator sums for DP40 are
$p+p^2+p^3$, $p^2+p^3+p^4$, and $p^3+p^4+p^5$, divided by $p^6-1$.
Cancelling gives the displayed triple; all sharp sequences follow because DP40 allowed arbitrary units. DP43 has codimension $6-3=3$ and supplies the exact affine loss. The mean and last-direction excess follow by the DP40 height-six algebra.
\end{proof}

This removes the conditional rank-one-display hypothesis for this concrete Jacobian at these primes. It leaves its nonalternating CM type unchanged and does not identify its equal-direction mean with a global theta cost.

